% GENERATED FILE. Edit canonical lessons, then run npm run generate:books.
% canonical-source-hash: fnv1a64:b92860e25ba563fc
% canonical-lessons: TE-C250-araceyyi, TE-C250-madama, TE-C250-cilamanda, TE-C250-toda, TE-C250-nadumu

\chapter{Palm of the hand, Heel, Ankle, Thigh, Waist}
\label{ch:te-palm-of-the-hand-heel-ankle-thigh-waist}
\hlchaptermodality{250}

\begin{chapteropening}
\textbf{By the end of this chapter:} \emph{I can say the palm of the hand, the heel, the ankle, the thigh and the waist.}

Complete the last lesson of chapter 250: I can say the palm of the hand, the heel, the ankle, the thigh and the waist.
\end{chapteropening}

\section[\texorpdfstring{\te{అరచెయ్యి} (araceyyi)}{అరచెయ్యి (araceyyi)}]{\te{అరచెయ్యి} (araceyyi) — the palm of the hand}

\label{lesson:TE-C250-araceyyi}



\begin{quote}
Before the new one: say the Telugu for to stretch out, then the Telugu for to stack.
\end{quote}

\subsection*{You'll want to know: \te{అరచెయ్యి}}
\textbf{\te{అరచెయ్యి}} — \emph{araceyyi} — \textquotedblleft{}the palm of the hand\textquotedblright{}.

\begin{grammarlens}[title={What you've built}]
One more word for this chapter.
\end{grammarlens}

\subsection*{Your turn}
\begin{itemize}
\raggedright
  \item \emph{Say it:} \emph{araceyyi}
  \item \emph{Say it:} \emph{araceyyi}, once more
  \item \emph{Say it:} say \emph{araceyyi}
  \item \emph{Recall:} say the Telugu for to stretch out, then the Telugu for to stack, then say \emph{araceyyi} again
\end{itemize}

\begin{tcolorbox}[breakable,colback=teal!4,colframe=teal!35!black,title={Before you move on}]
What does \te{అరచెయ్యి} mean? (\textquotedblleft{}The palm of the hand\textquotedblright{}.) Say it once more.
\end{tcolorbox}
\section[\texorpdfstring{\te{మడమ} (maḍama)}{మడమ (maḍama)}]{\te{మడమ} (maḍama) — the heel}

\label{lesson:TE-C250-madama}



\begin{quote}
Before the new one: say the Telugu for to stack, then the Telugu for the palm of the hand.
\end{quote}

\subsection*{You'll want to know: \te{మడమ}}
\textbf{\te{మడమ}} — \emph{maḍama} — \textquotedblleft{}the heel\textquotedblright{}.

\begin{grammarlens}[title={What you've built}]
One more word for this chapter.
\end{grammarlens}

\subsection*{Your turn}
\begin{itemize}
\raggedright
  \item \emph{Say it:} \emph{maḍama}
  \item \emph{Say it:} \emph{maḍama}, once more
  \item \emph{Say it:} say \emph{maḍama}
  \item \emph{Recall:} say the Telugu for to stack, then the Telugu for the palm of the hand, then say \emph{maḍama} again
\end{itemize}

\begin{tcolorbox}[breakable,colback=teal!4,colframe=teal!35!black,title={Before you move on}]
What does \te{మడమ} mean? (\textquotedblleft{}The heel\textquotedblright{}.) Say it once more.
\end{tcolorbox}
\section[\texorpdfstring{\te{చీలమండ} (cīlamaṇḍa)}{చీలమండ (cīlamaṇḍa)}]{\te{చీలమండ} (cīlamaṇḍa) — the ankle}

\label{lesson:TE-C250-cilamanda}



\begin{quote}
Before the new one: say the Telugu for the palm of the hand, then the Telugu for the heel.
\end{quote}

\subsection*{You'll want to know: \te{చీలమండ}}
\textbf{\te{చీలమండ}} — \emph{cīlamaṇḍa} — \textquotedblleft{}the ankle\textquotedblright{}.

\begin{grammarlens}[title={What you've built}]
One more word for this chapter.
\end{grammarlens}

\subsection*{Your turn}
\begin{itemize}
\raggedright
  \item \emph{Say it:} \emph{cīlamaṇḍa}
  \item \emph{Say it:} \emph{cīlamaṇḍa}, once more
  \item \emph{Say it:} say \emph{cīlamaṇḍa}
  \item \emph{Recall:} say the Telugu for the palm of the hand, then the Telugu for the heel, then say \emph{cīlamaṇḍa} again
\end{itemize}

\begin{tcolorbox}[breakable,colback=teal!4,colframe=teal!35!black,title={Before you move on}]
What does \te{చీలమండ} mean? (\textquotedblleft{}The ankle\textquotedblright{}.) Say it once more.
\end{tcolorbox}
\section[\texorpdfstring{\te{తొడ} (toḍa)}{తొడ (toḍa)}]{\te{తొడ} (toḍa) — the thigh}

\label{lesson:TE-C250-toda}



\begin{quote}
Before the new one: say the Telugu for the heel, then the Telugu for the ankle.
\end{quote}

\subsection*{You'll want to know: \te{తొడ}}
\textbf{\te{తొడ}} — \emph{toḍa} — \textquotedblleft{}the thigh\textquotedblright{}.

\begin{grammarlens}[title={What you've built}]
One more word for this chapter.
\end{grammarlens}

\subsection*{Your turn}
\begin{itemize}
\raggedright
  \item \emph{Say it:} \emph{toḍa}
  \item \emph{Say it:} \emph{toḍa}, once more
  \item \emph{Say it:} say \emph{toḍa}
  \item \emph{Recall:} say the Telugu for the heel, then the Telugu for the ankle, then say \emph{toḍa} again
\end{itemize}

\begin{tcolorbox}[breakable,colback=teal!4,colframe=teal!35!black,title={Before you move on}]
What does \te{తొడ} mean? (\textquotedblleft{}The thigh\textquotedblright{}.) Say it once more.
\end{tcolorbox}
\section[\texorpdfstring{\te{నడుము} (naḍumu)}{నడుము (naḍumu)}]{\te{నడుము} (naḍumu) — the waist}

\label{lesson:TE-C250-nadumu}



\begin{quote}
Before the new one: say the Telugu for the ankle, then the Telugu for the thigh.
\end{quote}

\subsection*{You'll want to know: \te{నడుము}}
\textbf{\te{నడుము}} — \emph{naḍumu} — \textquotedblleft{}the waist\textquotedblright{}.

\begin{grammarlens}[title={What you've built}]
One more word for this chapter.
\end{grammarlens}

\subsection*{Your turn}
\begin{itemize}
\raggedright
  \item \emph{Say it:} \emph{naḍumu}
  \item \emph{Say it:} \emph{naḍumu}, once more
  \item \emph{Say it:} say \emph{naḍumu}
  \item \emph{Recall:} say the Telugu for the ankle, then the Telugu for the thigh, then say \emph{naḍumu} again
\end{itemize}

\begin{tcolorbox}[breakable,colback=teal!4,colframe=teal!35!black,title={Before you move on}]
What does \te{నడుము} mean? (\textquotedblleft{}The waist\textquotedblright{}.) Say it once more.
\end{tcolorbox}
